\scanpage{5}
- \mspage{2} -

de plus, on vérifie fibre par fibre que
\[
R^1p_*\mathcal O(ne)=0\quad\text{pour}\quad n>0.
\]
Il en résulte, par Riemann-Roch, que pour $n>0$, le faisceau $p_*\mathcal O(ne)$ est localement libre de rang $n$, et que sa formation commute à tout changement de base. La suite exacte de cohomologie fournit de plus une suite exacte
\diagram{\begin{tikzcd}[column sep=small]
0\arrow[r]&p_*\mathcal O(ne)\arrow[r]&p_*\mathcal O((n+1)e)\arrow[r]&\omega^{\otimes-(n+1)}\arrow[r]&0
\end{tikzcd}\quad $(n>0)$.}
Par ailleurs, $\mathcal O(e)$ étant localement libre de rang 1, on vérifie fibre par fibre, compte tenu de (SGA 1 X 1.2), que
\diagram{\begin{tikzcd}
\mathcal O_S\arrow[r,"\sim"]&p_*\mathcal O_E\arrow[r,"\sim"]&p_*\mathcal O(e)
\end{tikzcd}.}
En d'autres termes, les $p_*\mathcal O(me)$ pour $1\leq m\leq n$ définissent une filtration de $p_*\mathcal O(ne)$, de gradué associé
\[
(1.1)\qquad \mathrm{Gr}\ p_*\mathcal O(ne)=\sum_{\substack{i=0\\i\ne1}}^n\omega^{\otimes-i}.
\]

Le morphisme $p$ est d'intersection complète relative ; il en résulte que $Rp^!\mathcal O_S\in D^+(E)$ \inserted{\unsure{$\rho$}} a un seul groupe de cohomologie non nulle, $R^1p^!\mathcal O_S$, et que ce dernier est un faisceau inversible sur $E$ ; on l'appelle le faisceau des différentielles régulières sur $E$ et on le désigne par $\Omega^{\mathrm{reg}}_{E/S}$. Là où $p$ est lisse, on a $\Omega^{\mathrm{reg}}_{E/S}=\Omega^1_{E/S}$. La dualité montre que $R^1p_*\Omega^{\mathrm{reg}}_{E/S}$ est isomorphe à $\mathcal O$, $p_*\Omega^{\unsure{\mathrm{rog}|\mathrm{reg}}}_{E/S}$ est donc\inserted{\unsure{$\rho$}} localement libre, de formation compatible à tout changement de base et, raisonnant fibre par fibre, on en déduit que
\margin{\uline{réf}}
\margin{\uline{réf}}
\margin{4}
